\section{Результаты моделирования}

\subsection{Параметры модели}

Моделирование проведено для радиального подшипника скольжения геометрии~\varGeomLetter{} ($R = \varRmm$~мм, $c = \varCmkm$~мкм, $L = \varLmm$~мм), поверхность которого содержит \varDimpleTypeGenitive{} микрорельефом. В таблице~\ref{tab:input_data} систематизированы исходные данные.

\begin{table}[H]
\small
\caption{Исходные данные для расчёта}
\label{tab:input_data}
\begin{tabularx}{\textwidth}{|X|c|c|}
\hline
\centering\textbf{Параметр} & \textbf{Значение} & \textbf{Размерность} \\
\hline
Радиус вала $R$ & \varRmm & мм \\
\hline
Радиальный зазор $c$ & \varCmkm & мкм \\
\hline
Длина подшипника $L$ & \varLmm & мм \\
\hline
Относительная длина $L/D$ & \varLD & -- \\
\hline
Полуось по оси $Z$, $a$ & \varDimpleAmm & мм \\
\hline
Полуось по оси $\varphi$, $b$ & \varDimpleBmm & мм \\
\hline
Глубина углубления $h_p$ & \varDimpleHpmkm & мкм \\
\hline
Относительная глубина $h_p/c$ & \varDimpleHpRatio & -- \\
\hline
Число углублений по окружности $N_\varphi$ & \varNphi & шт. \\
\hline
Число рядов по длине $N_Z$ & \varNZ & шт. \\
\hline
Область текстурирования & \varPhiRange & -- \\
\hline
Частота вращения $n$ & \varN & об/мин \\
\hline
Вязкость смазки $\eta$ & \varEta & Па$\cdot$с \\
\hline
Диапазон эксцентриситета $\varepsilon$ & 0{,}05-0{,}80 & -- \\
\hline
Размер сетки & $500 \times 500$ & -- \\
\hline
Параметр релаксации $\omega$ & 1{,}5 & -- \\
\hline
Критерий сходимости $\delta$ & $< 10^{-5}$ & -- \\
\hline
Макс. число итераций & 50000 & -- \\
\hline
\end{tabularx}
\end{table}

При заданной частоте вращения линейная скорость вала равна $U = \pi D n / 60 = \pi \cdot 0{,}100 \cdot 2980 / 60 \approx 15{,}60$~м/с. Расчёт проведён для 15 значений эксцентриситета в указанном диапазоне.

\subsection{Распределение зазора и давления}

На рисунке~\ref{fig:pressure_2d} показано распределение безразмерного давления $P(\varphi)$ в среднем сечении ($Z = 0$) при $\varepsilon = 0{,}6$ для гладкого и текстурированного подшипников.

\begin{figure}[H]
\centering
\includegraphics[width=\textwidth]{figures/pressure_2d.png}
\caption{Распределение давления $P(\varphi)$ при $Z=0$ для гладкого и текстурированного подшипников ($\varepsilon = 0{,}6$)}
\label{fig:pressure_2d}
\end{figure}

В случае гладкого подшипника давление монотонно возрастает в сужающейся части зазора и достигает единственного максимума вблизи $\varphi \approx \pi$. В расширяющейся части зазора давление обращается в нуль вследствие кавитации, а зона положительного давления охватывает приблизительно половину окружности подшипника. Введение микрорельефа приводит к возникновению локальных осцилляций давления на каждом углублении: входная кромка углубления формирует локальное повышение давления, а выходная -- его снижение. Суммарный эффект проявляется в увеличении общего уровня давления, что обусловливает рост несущей способности.

Профиль зазора $H(\varphi)$ в среднем сечении при $\varepsilon = 0{,}6$ приведён на рисунке~\ref{fig:clearance_2d}.

\begin{figure}[H]
\centering
\includegraphics[width=\textwidth]{figures/clearance_2d.png}
\caption{Распределение зазора $H(\varphi)$ при $Z=0$ для подшипника с \varDimpleTypeGenitive{} углублениями ($\varepsilon = 0{,}6$)}
\label{fig:clearance_2d}
\end{figure}

Зазор изменяется по закону $H = 1 + \varepsilon\cos\varphi$, а эллипсоидальные углубления проявляются в виде локальных увеличений зазора. Безразмерная глубина каждого углубления составляет $\Delta H_{\max} = h_p/c = 0{,}1$, что отражает меньшее соотношение глубины к зазору по сравнению с подшипниками меньшего типоразмера.

На рисунках~\ref{fig:clearance_3d} и~\ref{fig:pressure_3d} представлены трёхмерные распределения зазора и давления при $\varepsilon = 0{,}6$.

\begin{figure}[H]
\centering
\includegraphics[width=\textwidth]{figures/clearance_3d.png}
\caption{Трёхмерное распределение зазора $H(\varphi, Z)$ ($\varepsilon = 0{,}6$)}
\label{fig:clearance_3d}
\end{figure}

Регулярная структура микрорельефа отчётливо прослеживается на рисунке~\ref{fig:clearance_3d}: 8 углублений по окружности, расположенных в 11 рядах по длине, сосредоточены в зоне $90^\circ - 270^\circ$.

\begin{figure}[H]
\centering
\includegraphics[width=\textwidth]{figures/pressure_3d.png}
\caption{Трёхмерное распределение давления $P(\varphi, Z)$ ($\varepsilon = 0{,}6$)}
\label{fig:pressure_3d}
\end{figure}

Помимо основного максимума давления, формируемого макрогеометрией клинового зазора, на рисунке~\ref{fig:pressure_3d} видны локальные возмущения, порождённые углублениями. Сужение зазора на входной кромке углубления вызывает дополнительный подъём давления -- микрогидродинамический эффект. На выходной кромке давление снижается, но кавитационное ограничение ($P \geq 0$) обеспечивает положительный суммарный вклад каждого углубления в несущую способность.

Из анализа трёхмерных полей следуют характерные закономерности. Максимум давления достигается в среднем сечении ($Z = 0$) и спадает к торцам подшипника ($Z = \pm 1$), где давление обращается в нуль из-за торцевых утечек. Эффект текстурирования наиболее выражен в центральной области, где фоновое давление велико. Для периферийных рядов углублений, расположенных вблизи торцов, влияние микрорельефа ослабевает.

\subsection{Сравнительный анализ интегральных показателей}

В таблицах~\ref{tab:smooth_results} и~\ref{tab:textured_results} приведены интегральные характеристики подшипника для выбранных значений эксцентриситета.

\begin{table}[H]
\small
\centering
\caption{Характеристики гладкого подшипника}
\label{tab:smooth_results}
\begin{tabular}{|c|c|c|c|}
\hline
$\boldsymbol{\varepsilon}$ & $\boldsymbol{F}$\textbf{, Н} & $\boldsymbol{\mu}$ & $\boldsymbol{Q}$\textbf{, л/с} \\
\hline
0{,}05 & 551{,}5 & 0{,}2245 & 0{,}6852 \\
\hline
0{,}16 & 1775{,}8 & 0{,}0707 & 0{,}6843 \\
\hline
0{,}26 & 3141{,}8 & 0{,}0412 & 0{,}6825 \\
\hline
0{,}37 & 4805{,}2 & 0{,}0282 & 0{,}6801 \\
\hline
0{,}48 & 6984{,}6 & 0{,}0207 & 0{,}6772 \\
\hline
0{,}59 & 10124{,}4 & 0{,}0157 & 0{,}6744 \\
\hline
0{,}69 & 15347{,}9 & 0{,}0118 & 0{,}6717 \\
\hline
0{,}80 & 25964{,}3 & 0{,}0085 & 0{,}6698 \\
\hline
\end{tabular}
\end{table}

\begin{table}[H]
\small
\centering
\caption{Характеристики подшипника с \varDimpleTypeGenitive{} углублениями}
\label{tab:textured_results}
\begin{tabular}{|c|c|c|c|}
\hline
$\boldsymbol{\varepsilon}$ & $\boldsymbol{F}$\textbf{, Н} & $\boldsymbol{\mu}$ & $\boldsymbol{Q}$\textbf{, л/с} \\
\hline
0{,}05 & 609{,}3 & 0{,}2023 & 0{,}6888 \\
\hline
0{,}16 & 1857{,}8 & 0{,}0672 & 0{,}6878 \\
\hline
0{,}26 & 3249{,}7 & 0{,}0396 & 0{,}6861 \\
\hline
0{,}37 & 4957{,}0 & 0{,}0271 & 0{,}6836 \\
\hline
0{,}48 & 7222{,}0 & 0{,}0199 & 0{,}6808 \\
\hline
0{,}59 & 10609{,}3 & 0{,}0148 & 0{,}6778 \\
\hline
0{,}69 & 16417{,}7 & 0{,}0109 & 0{,}6751 \\
\hline
0{,}80 & 29831{,}7 & 0{,}0073 & 0{,}6729 \\
\hline
\end{tabular}
\end{table}

С увеличением эксцентриситета нагрузочная способность $F$ возрастает нелинейно для обоих вариантов подшипника, причём текстурированный подшипник демонстрирует более высокие значения $F$ во всём исследованном диапазоне. Наиболее значительная разница наблюдается при больших эксцентриситетах: при $\varepsilon = 0{,}80$ нагрузочная способность текстурированного подшипника достигает 29831{,}7~Н, что на 14{,}9\% выше значения для гладкого аналога (25964{,}3~Н).

Коэффициент трения $\mu$ убывает монотонно с ростом эксцентриситета, и для подшипника с микрорельефом значения $\mu$ оказываются ниже во всём диапазоне. Расход смазки $Q$ остаётся практически постоянным и для текстурированного подшипника несколько превышает расход гладкого за счёт увеличения эффективного объёма смазочного слоя.

Количественная оценка эффективности текстурирования при $\varepsilon = 0{,}6$ приведена в таблице~\ref{tab:comparison}.

\begin{table}[H]
\small
\centering
\caption{Сравнение характеристик при $\varepsilon = 0{,}6$}
\label{tab:comparison}
\begin{tabular}{|l|c|c|c|}
\hline
\textbf{Характеристика} & \textbf{Гладкий} & \textbf{С углублениями} & $\boldsymbol{\Delta}$\textbf{, \%} \\
\hline
Нагрузка $F$, Н & 10685{,}6 & 11175{,}6 & $+4{,}59$ \\
\hline
Коэфф. трения $\mu$ & 0{,}0151 & 0{,}0143 & $-5{,}21$ \\
\hline
Расход $Q$, л/с & 0{,}6739 & 0{,}6774 & $+0{,}51$ \\
\hline
\end{tabular}
\end{table}

Эллипсоидальные углубления обеспечивают прирост несущей способности на 4{,}59\% благодаря микрогидродинамическому эффекту на кромках углублений. Снижение коэффициента трения составляет 5{,}21\%, что объясняется увеличением средней толщины смазочного слоя в зонах текстурирования. Расход смазки увеличивается лишь на 0{,}51\%, что не создаёт дополнительных требований к системе маслоснабжения и одновременно способствует улучшению отвода тепла из зоны трения.

Зависимости интегральных характеристик от эксцентриситета для обоих вариантов подшипника представлены на рисунках~\ref{fig:load_vs_eps}-\ref{fig:flow_vs_eps}.

\begin{figure}[H]
\centering
\includegraphics[width=\textwidth]{figures/load_vs_epsilon.png}
\caption{Нагрузочная способность $F(\varepsilon)$ для гладкого и текстурированного подшипников}
\label{fig:load_vs_eps}
\end{figure}

Во всём диапазоне эксцентриситетов кривая $F(\varepsilon)$ текстурированного подшипника располагается выше кривой гладкого. Различие между кривыми нарастает с увеличением $\varepsilon$: при $\varepsilon = 0{,}05$ абсолютная разница составляет около 58~Н, тогда как при $\varepsilon = 0{,}80$ она достигает порядка 3870~Н. Такое поведение объясняется степенной зависимостью давления от зазора ($P \sim h^{-3}$), в силу чего при малых зазорах влияние текстурирования усиливается.

\begin{figure}[H]
\centering
\includegraphics[width=\textwidth]{figures/friction_vs_epsilon.png}
\caption{Коэффициент трения $\mu(\varepsilon)$ для гладкого и текстурированного подшипников}
\label{fig:friction_vs_eps}
\end{figure}

Для обоих вариантов коэффициент трения монотонно убывает и при $\varepsilon \to 1$ асимптотически стремится к нулю. Текстурированный подшипник показывает более низкие значения $\mu$ при всех эксцентриситетах. Наибольшая абсолютная разница зафиксирована при малых $\varepsilon$: при $\varepsilon = 0{,}05$ она составляет $\Delta\mu = 0{,}022$, тогда как при $\varepsilon = 0{,}80$ -- $\Delta\mu = 0{,}001$. Тем не менее в относительном выражении снижение трения сохраняется на сопоставимом уровне во всём диапазоне.

\begin{figure}[H]
\centering
\includegraphics[width=\textwidth]{figures/flow_vs_epsilon.png}
\caption{Расход смазки $Q(\varepsilon)$ для гладкого и текстурированного подшипников}
\label{fig:flow_vs_eps}
\end{figure}

Расход смазки текстурированного подшипника несколько превышает расход гладкого, однако различие не превышает 0{,}5\% и практически не зависит от уровня эксцентриситета.

Результаты свидетельствуют о том, что текстурирование поверхности подшипника геометрии~B позволяет одновременно повысить нагрузочную способность и уменьшить трение при минимальном влиянии на расход смазки. Эффект текстурирования возрастает при увеличении эксцентриситета, что имеет важное значение для тяжелонагруженных узлов оборудования.

\subsection*{Выводы по главе 3}
\phantomsection
\addcontentsline{toc}{subsection}{Выводы по главе 3}

Проведён расчёт характеристик радиального подшипника скольжения геометрии~B с эллипсоидальным микрорельефом для 15 значений эксцентриситета. Построены поля давления и зазора, определены зависимости нагрузочной способности, коэффициента трения и расхода смазки от эксцентриситета. Сравнение характеристик гладкого и текстурированного подшипников при $\varepsilon = 0{,}6$ выявило прирост нагрузочной способности на $4{,}59$\%, снижение коэффициента трения на $5{,}21$\% и увеличение расхода на $0{,}51$\%, что подтверждает положительное влияние эллипсоидального текстурирования.
